\documentclass[11pt]{article}

\usepackage{amsmath,amssymb}
\usepackage{xurl}
\usepackage[hidelinks]{hyperref}
\setlength{\emergencystretch}{2em}

\input{command}

% Local notation for this note.
\DeclareMathOperator{\poly}{poly}
\newcommand{\Id}{\mathrm{Id}}
\renewcommand{\ghissue}[1]{%
  \href{https://github.com/Dengnifer/MIPStarRE-A/issues/#1}{GitHub issue \##1}}

\title{The Line Marginal in the Second Sub-Line Claim}
\author{}
\date{2026-09-05}

\begin{document}
\maketitle

\section{At a glance}

\begin{itemize}
  \item \textbf{Difficulty: easy.}  The second of the three scalar claims
  in the proof of the combined-lines lemma of the Pauli basis analysis
  removes the trailing $X$-point factor at the cost
  $m\,\delta_{\mathrm{Line}}^{1/2}$.  Its printed proof identifies the
  $X$-marginal of the paired-line measurement with the evaluation class of
  the $X$-line measurement through the sandwich form of the paired-line
  measurement.  That form is the construction inside the proof of the
  two-basis line lemma; the conclusion of that lemma records only the
  consistency of the evaluated pairs with the joint point measurement, and
  a measurement satisfying the conclusion need not have the $X$-line
  measurement as its marginal.  The claim as printed therefore does not
  follow from the objects the lemmas provide, and for an arbitrary witness
  of the conclusion it is false.  The correction proves the claim with the
  error $\sqrt{m}\,(\delta_P^{1/4} + \delta_Q^{1/4})$ from the consistency
  conclusion alone; the sequel is unaffected.
  \item \textbf{Key Mathlib inputs:} the continuous functional calculus
  square root of a positive matrix, for the Cauchy--Schwarz inequality of a
  positive operator weight, together with finite averages and the
  state-dependent distance calculus of the project (blueprint
  Chapter~12).
  \item \textbf{Estimated weight:} a \textbf{statement-level error-form
  correction} of one claim, plus about \textbf{four lemma interfaces}: the
  weighted Cauchy--Schwarz estimate in deficit form, the one-point
  deficits of the evaluated paired-line measurement over a product of
  restricted line-point laws, their transport to the sub-line law, and
  the numerical comparison of the two error forms; about \textbf{one
  thousand lines} of formal proof in total, now complete.
  \item \textbf{Mathlib/project split:} essentially \textbf{100\%
  project-local}, the only Mathlib surface being the square root of a
  positive matrix and elementary real inequalities.
  \item \textbf{Key inputs:} the combined point and paired-line
  measurements and the sub-line distribution of the Pauli basis analysis
  in \cite{Ji2020MIPStar}.  The formalization is tracked in
  \ghissue{118}; the error term of the enclosing lemma is the subject of
  \path{docs/paper-gaps/qpbt_combined-lines-error-term.tex}.
\end{itemize}

\section{Key theorem forms}

Throughout, $(q,m,d)$ is an admissible parameter tuple of the Pauli basis
test, $\ket{\hat\psi}$ is the expanded state of the projective strategy
under analysis, succeeding with probability at least $1-\eps$, and all
approximations are with respect to $\ket{\hat\psi}$.  For lines
$\ell_X, \ell_Z$ in $\F_q^m$, $\{T^{\ell_X,\ell_Z}_{f_X,f_Z}\}$ is the
paired-line measurement of the two-basis line lemma, with outcomes pairs
of polynomials of degree at most $md$ on the two lines, and
$\{\hat Q^{x,z}_{a,b}\}$ is the combined point measurement; their
consistency errors are $\delta_P = \poly(\eps, md/q)$ and
$\delta_Q = \poly(\eps)$, and $\delta_{\mathrm{Line}} = \poly(\eps)$ is the
line--point consistency error of the expanded line measurements
$\{\hat M^{(\mathrm{Line},W),\ell}_f\}$.  The sub-line distribution $D$ samples
a line $\ell$ of $\F_q^{2m+2}$ together with two lines $\ell_X, \ell_Z$ of
$\F_q^m$, and $(x,z,\alpha,\beta)$ is a uniformly random point of
$\ell$.%
\footnote{In the mirror of the source consulted here,
\path{references/qpbt-paper/14_analysis_of_the_pauli_basis_test.tex}: the
two-basis line lemma is at lines~882--894 (label \path{lem:qld-xz-lines}),
with the sandwich definition of $T$ inside its proof; the sub-line lemma at
lines~1063--1069 (label \path{lem:qld-sublines}); the second claim at
lines~1168--1201 (label \path{claim:17-2}), whose final lines use the
sandwich form and the projectivity of the line measurements; and the third
claim at lines~1204--1239 (label \path{claim:17-3}).}
The claim under discussion asserts
\begin{align*}
  \E_{(\ell,\ell_X,\ell_Z) \sim D}\ \E_{(x,z,\alpha,\beta) \in \ell}
  &\sum_{f_X,f_Z} \bra{\hat\psi}
  \big(T^{\ell_X,\ell_Z}_{f_X,f_Z}\big)_{\mathsf{A}\mathsf{A}'} \ot
  \big(\hat{M}^{(\mathrm{Point},Z),z}_{f_Z(z)}
  \hat{M}^{(\mathrm{Point},X),x}_{f_X(x)}\big)_{\mathsf{B}\mathsf{A}''}
  \ket{\hat\psi} \\
  \approx_{\delta}\
  \E_{(\ell,\ell_X,\ell_Z) \sim D}\ \E_{(x,z,\alpha,\beta) \in \ell}
  &\sum_{f_X,f_Z} \bra{\hat\psi}
  \big(T^{\ell_X,\ell_Z}_{f_X,f_Z}\big)_{\mathsf{A}\mathsf{A}'} \ot
  \big(\hat{M}^{(\mathrm{Point},Z),z}_{f_Z(z)}\big)_{\mathsf{B}\mathsf{A}''}
  \ket{\hat\psi}
\end{align*}
with the following error forms.
\begin{enumerate}
  \item \textbf{Source form.}  $\delta = m\,\delta_{\mathrm{Line}}^{1/2}$.
  Its derivation uses the identity
  $\sum_{f_Z} T^{\ell_X,\ell_Z}_{f_X,f_Z} =
  \hat{M}^{(\mathrm{Line},X),\ell_X}_{f_X}$, which holds for the sandwich
  $T^{\ell_X,\ell_Z}_{f_X,f_Z} = \hat{M}^{(\mathrm{Line},X),\ell_X}_{f_X}
  \hat{M}^{(\mathrm{Line},Z),\ell_Z}_{f_Z}
  \hat{M}^{(\mathrm{Line},X),\ell_X}_{f_X}$ but is not part of the
  conclusion of the two-basis line lemma.
  \item \textbf{Corrected form (adopted).}
  $\delta = \sqrt{m}\,(\delta_P^{1/4} + \delta_Q^{1/4})$, with the explicit
  constant $2$.  It follows from the consistency conclusion of the
  two-basis line lemma, the two ordered displays of the combined point
  lemma, and Property~2 of the sub-line lemma.
  \item \textbf{Consumer requirement.}  The claim is chained with the
  first and third claims to bound the consistency defect of the
  combined-lines lemma by
  $O(m(\eps^{1/4} + \delta_P^{1/4} + \delta_Q^{1/4} +
  \delta_{\mathrm{Line}}^{1/2}))$; the corrected form is dominated by that
  bound, so the chain and its consumer are unchanged.
\end{enumerate}

\section{The source assertion and why the objects do not support it}

The printed proof of the second claim applies the Cauchy--Schwarz
inequality with both factors weighted by
$\sqrt{T^{\ell_X,\ell_Z}_{f_X,f_Z}}$ and reduces the claim to bounding the
deficit
\begin{equation*}
  1 - \E_{(\ell,\ell_X,\ell_Z) \sim D}\ \E_{(x,z,\alpha,\beta) \in \ell}
  \sum_a \bra{\hat\psi}
  \big(T^{\ell_X,\ell_Z,x}_{a}\big)_{\mathsf{A}\mathsf{A}'} \ot
  \big(\hat{M}^{(\mathrm{Point},X),x}_{a}\big)_{\mathsf{B}\mathsf{A}''}
  \ket{\hat\psi}\,,
  \qquad
  T^{\ell_X,\ell_Z,x}_{a} = \sum_{f_X,f_Z \,:\, f_X(x)=a}
  T^{\ell_X,\ell_Z}_{f_X,f_Z}\,.
\end{equation*}
It then asserts, ``using the definition'' of $T$ as a sandwich ``and the
fact that'' the $X$-line measurement ``is projective'', that
$T^{\ell_X,\ell_Z,x}_{a} =
\hat{M}^{(\mathrm{Line},X),\ell_X}_{[\mathrm{eval}_x(\cdot)=a]}$, after
which Property~2 of the sub-line lemma and item~3 of the line consistency
lemma bound the deficit by $O(m^2\delta_{\mathrm{Line}})$.

The identity is correct for the sandwich, since
$\sum_{f_Z} \hat{M}^{X}_{f_X} \hat{M}^{Z}_{f_Z} \hat{M}^{X}_{f_X}
= \hat{M}^{X}_{f_X} \hat{M}^{X}_{f_X} = \hat{M}^{X}_{f_X}$.  But the
two-basis line lemma is stated existentially: it provides, for each pair
of lines, a measurement $\{T^{\ell_X,\ell_Z}_{f_X,f_Z}\}$ whose evaluated
pairs are consistent with the joint point measurement, on average over two
independent line-point pairs, together with the degree condition on
axis-parallel lines.  The sandwich form is the witness exhibited in its
proof, not a property carried by its conclusion, and the sub-line lemma and
the three claims are stated for the measurement of the two-basis line
lemma, that is, for a measurement satisfying that conclusion.  For such a
measurement the marginal identity can fail: with a strategy of small
$\eps$, the family
$T^{\ell_X,\ell_Z}_{f_X,f_Z} = [f_X = 0]\,
\hat{M}^{(\mathrm{Line},Z),\ell_Z}_{f_Z}$ is a measurement whose
consistency defect is trivially at most $\delta_P = 1$, while the
difference of the two sides of the claim is, up to an error of order
$\delta_Q^{1/2}$, the mass $1 - \E_x \bra{\hat\psi}
\hat{M}^{(\mathrm{Point},X),x}_{0} \ket{\hat\psi}$ of the nonzero
$X$-point answers, which does not tend to $0$ with $\eps$.  The claim with
the source error is therefore false for some witnesses of the two-basis
line lemma, although true for the sandwich witness.

\section{Correction and alternative repairs}

\paragraph{Corrected error form (adopted in the blueprint and the
formalization).}
The deficit is bounded from the consistency conclusion alone.  The
quantity averaged depends on the sample only through $(\ell_X, \ell_Z,
x)$, so by Property~2 of the sub-line lemma the average is a mixture of
averages over products $D_{\mathrm{v},i} \times D_{\mathrm{v},j}$ of two
restricted line-point laws, in which the point $z'$ of the second factor is
fresh and unused; and $T^{\ell_X,\ell_Z,x}_a = \sum_b
T^{\ell_X,\ell_Z}_{[\mathrm{eval}_x(\cdot)=a,\ \mathrm{eval}_{z'}(\cdot)=b]}$
for every $z'$.  Over one product law, with $T_{a,b}$ the evaluation
classes at $(x, z')$, $\hat Q_{a,b}$ the completed joint point measurement
at $(x,z')$ placed on $\mathsf{B}\mathsf{A}''$, and $\hat M^X_a$, $\hat
M^Z_b$ the two point measurements there, three Cauchy--Schwarz estimates
weighted by $T_{a,b}$ give
\begin{equation*}
  \E \sum_{a,b} \bra{\hat\psi} T_{a,b} \ot \hat M^X_a \ket{\hat\psi}
  \ \ge\ \E \sum_{a,b} \bra{\hat\psi} T_{a,b} \ot \hat Q_{a,b}
  \ket{\hat\psi} - \eta^{1/2} - 2\zeta^{1/2}
  \ \ge\ 1 - 2\eta^{1/2} - 2\zeta^{1/2}\,,
\end{equation*}
where $\eta \le 4m^2\delta_P$ is the consistency defect of the evaluated
pairs with $\hat Q$ over the product law (item~2 of the restricted-mixture
lemma) and $\zeta = 4\delta_Q$ bounds the squared distance of $\hat
Q_{a,b}$ from each ordered product of $\hat M^X_a$ and $\hat M^Z_b$ on one
placement: the first estimate inserts $\hat Q_{a,b}$ before the projection
$\hat M^X_a$, the second replaces $\hat Q_{a,b}$ by
$\hat M^Z_b \hat M^X_a$, which absorbs the trailing factor, and the third
compares the overlaps of $T_{a,b}$ with the ordered product and with $\hat
Q_{a,b}$.  The deficit is thus at most $4m\,\delta_P^{1/2} +
4\,\delta_Q^{1/2}$, and its square root at most
$2\sqrt{m}\,(\delta_P^{1/4} + \delta_Q^{1/4})$.  The same argument, with the
roles of the two points and of the two ordered products exchanged, bounds
the deficit appearing in the third claim, and shows in passing that the
Cauchy--Schwarz step of the third claim is vacuous: its right-hand side is
the square root of its left-hand side, so the third claim is the deficit
bound itself, without the $\eps^{1/4}$ term the source obtains from the
self-consistency of the points measurements.

\paragraph{Restoring the source error form.}
Strengthen the conclusion of the two-basis line lemma by the marginal
identity $\sum_{f_Z} T^{\ell_X,\ell_Z}_{f_X,f_Z} =
\hat{M}^{(\mathrm{Line},X),\ell_X}_{f_X}$ for all $\ell_X, \ell_Z, f_X$.
The sandwich witness satisfies it, so the lemma remains true, and the
printed proof of the second claim then goes through with the error
$(2m\,\delta_{\mathrm{Line}})^{1/2}$, item~3 of the line consistency lemma
being inflated by $2m$ on one restricted law.  This is a change to a source
statement and to the carrier of its conclusion in the formalization; it is
recorded here as the alternative and is not adopted.  Note that the
$Z$-marginal of the sandwich is not the $Z$-line measurement, so the
analogous identity is not available for the third claim, whose source proof
does not use it.

\section{Comparison with the blueprint}

The blueprint \cite{MIPStarREBlueprint} states the second claim with the
corrected error form, keeps the source form in a comment and in a remark
that cites this note, and proves it along the route above; the third claim
keeps its source statement, with the source's $\eps^{1/4}$ term retained
for fidelity, and its proof is replaced by the deficit bound.  The
Cauchy--Schwarz estimate in deficit form and the one-point deficits are
formalization-support nodes.  The chained bound in the proof of the
combined-lines lemma is unchanged, the corrected error being dominated by
the printed one.%
\footnote{The blueprint statements are \path{lem:claim-17-2},
\path{lem:claim-17-3}, \path{lem:overlap-gap-deficit},
\path{lem:qld-sublines-one-point-deficit}, and the remark
\path{rem:claim-17-2-line-marginal}, all in
\path{blueprint/src/chapter/ch15_qpbt_combining.tex}.}

The three claims are the declarations
\leanid{MIPStarRE.QPBT.subline\_replace\_by\_ordered\_product},
\leanid{MIPStarRE.QPBT.subline\_remove\_X\_factor}, and
\leanid{MIPStarRE.QPBT.subline\_Z\_term\_near\_one}; all three are proved,
the second with the corrected error form and the explicit constant $2$.
The paired-line measurement is the field of the record
\leanid{MIPStarRE.QPBT.CombinedLinesWitness}, which carries the
consistency conclusion and the degree condition only; the source error
form would require the marginal identity as an additional field of that
record and a corresponding strengthening of \path{lem:qld-xz-lines}.  The
Cauchy--Schwarz estimates in deficit form are in
\path{MIPStarRE/QPBT/Combining/OverlapGap.lean}, the one-point deficits in
\path{MIPStarRE/QPBT/Combining/EvalDeficit.lean}, and their transport to
the sub-line law in \path{MIPStarRE/QPBT/Combining/SubLineDeficit.lean}.

\section{Conclusion}

There is a proof-level gap in the second claim: its printed argument uses
the sandwich form of the paired-line measurement, which is the witness
constructed in the proof of the two-basis line lemma and not a property of
that lemma's conclusion, and the claim with the source error
$m\,\delta_{\mathrm{Line}}^{1/2}$ fails for other witnesses.  The claim
holds with the error $\sqrt{m}\,(\delta_P^{1/4} + \delta_Q^{1/4})$ from the
consistency conclusion alone, by the same Cauchy--Schwarz reduction
followed by a three-step comparison with the joint point measurement; the
third claim admits the same treatment, without its $\eps^{1/4}$ term.  The
sequel is unaffected.  The alternative repair, adding the marginal identity
to the conclusion of the two-basis line lemma, would restore the source
error and is a source-statement change left to the owner.

\bibliographystyle{alpha}
\bibliography{references}

\end{document}
